% W5i79_HardProblems.tex
% DFD 20-Agent Research Programme --- Iteration 79
% Hard Problems, Honest Negatives, and Open Questions

\documentclass[11pt,a4paper]{article}
\usepackage[utf8]{inputenc}
\usepackage{amsmath,amssymb}
\usepackage{booktabs}
\usepackage{geometry}
\usepackage{xcolor}
\usepackage{enumitem}
\geometry{margin=2.5cm}

\title{W5i79: Hard Problems and Honest Negatives\\[6pt]
\large Publication Risks $\cdot$ Community Reception $\cdot$ Referee Scenarios}
\author{DFD 20-Agent Research Programme}
\date{2026-03-24}

\begin{document}
\maketitle

\section{Hard Problem 1: PRL Rejection Risk}

The $\alpha$ PRL letter is the highest-stakes submission. Potential rejection scenarios:
\begin{enumerate}[nosep]
\item \textbf{``Numerology'' dismissal}: Referee argues that $\alpha^{-1} = 137.036\ldots$ could be coincidental curve-fitting. \textbf{Counter}: Zero free parameters. The formula has no adjustable constants. The inputs ($k=60$, $\Tr(Y^2)=10$) are topologically/group-theoretically determined. This is not fitting.

\item \textbf{``Unverified framework'' objection}: Referee argues DFD is not established physics and PRL should not publish speculative results. \textbf{Counter}: PRL publishes theoretical predictions that make testable predictions. DFD has 28 testable predictions, 5 imminent.

\item \textbf{``Too good to be true'' scepticism}: Sub-ppm agreement may trigger extraordinary scepticism. \textbf{Counter}: Present the honest negative (0.55 ppm theoretical uncertainty). The result is extraordinary but the methodology is transparent.

\item \textbf{Format rejection}: PRL may argue a 4-page letter cannot adequately present the derivation. \textbf{Counter}: Full derivation in supplemental material. Letter presents result + key steps + predictions.
\end{enumerate}

\textbf{Backup}: If PRL rejects, submit to Physical Review D (longer format, same society).

\textbf{Risk level}: MODERATE. PRL acceptance rate for theoretical physics is $\sim 25\%$. The result is exceptional but controversial.

\section{Hard Problem 2: Physics Reports Review Timeline}

Physics Reports average review time: 6--12 months. For a 248-page paper, potentially longer. Risks:
\begin{itemize}[nosep]
\item Reviewer fatigue: finding 2--3 referees willing to review 248 pages.
\item Scope: referees may argue the paper tries to do too much.
\item Delay: if review takes 12+ months, Euclid DR1 data may arrive during review, requiring revisions.
\end{itemize}

\textbf{Mitigation}: Cover letter emphasises modular structure (referees can focus on their expertise area). Offer to split into shorter papers if editor requests.

\section{Hard Problem 3: Community Reception --- ``Single Author'' Challenge}

DFD is primarily a single-author project. In modern physics:
\begin{itemize}[nosep]
\item Single-author papers are unusual in cosmology (typical: 3--20 authors).
\item May face implicit bias (``one person can't do all this'').
\item Collaboration engagement (i79 plan) is partially a response to this.
\end{itemize}

\textbf{Mitigation}: (a) Zenodo archive enables independent verification. (b) Conference presentations enable direct engagement. (c) Collaboration with external researchers adds credibility. (d) The mathematical content is self-verifying (proofs are proofs regardless of author count).

\section{Hard Problem 4: Simultaneous arXiv Posting}

Should all 8 papers be posted to arXiv simultaneously or sequentially?
\begin{itemize}[nosep]
\item \textbf{Simultaneous}: Maximum impact, establishes priority for all results at once. Risk: community overwhelmed, individual papers get less attention.
\item \textbf{Sequential}: Each paper gets dedicated attention. Risk: later papers may be scooped if ideas are public from v3.3.
\item \textbf{Recommended}: Phase 1 (PRL + v3.3) simultaneously. Then 1 paper per 2--3 weeks.
\end{itemize}

\section{Hard Problem 5: N-Body External Code Comparison}

The N-body paper (40\%) includes a comparison with MG-GADGET (Winther et al.\ 2019). But:
\begin{itemize}[nosep]
\item MG-GADGET implements Hu-Sawicki $f(R)$ and nDGP, not DFD specifically.
\item The comparison is with a ``Yukawa-coupled scalar'' model that has similar (but not identical) scale dependence.
\item Referees may object that this is not a true code comparison.
\end{itemize}

\textbf{Mitigation}: (a) State explicitly that the comparison is approximate. (b) Plan for a DFD-specific branch in a public N-body code (e.g., AREPO or RAMSES fork) in Phase 2. (c) The agreement at $< 0.5\%$ is a consistency check, not a validation.

\section{Open Questions from i79}

\begin{enumerate}[nosep]
\item Should v3.3 arXiv posting precede or coincide with Physics Reports submission?
\item Is there a COSMO 2026 abstract deadline in May that requires immediate action?
\item Should the N-body paper include a DFD implementation manual as an appendix (for reproducibility)?
\item What is the expected arXiv ``new submissions'' visibility for a flagship posting?
\item Should the collaboration outreach emails include a draft paper or just an abstract?
\end{enumerate}

\section{Summary}

5 hard problems. 0 RED risks. All are publication/community challenges, not scientific issues. The science is settled (97/3/0, adversarially audited). The remaining challenges are sociological: how to present extraordinary claims to a sceptical community in the most effective way. Programme integrity: \textbf{ABSOLUTE}.

\end{document}
